\section{Personnes commissaires}
\label{personnes-commissaires}

\subsection{Nomination et ratification}
\label{personnes-commissaires-nomination-et-ratification}

\begin{enumerate}
 \item Le CA, sur recommandation de l'exécutif national, peut nommer les personnes commissaires aux fins d'aider avec les portefeuilles respectifs de l'exécutif national.
\end{enumerate}

\subsection{Mandat}
\label{personnes-commissaires-mandat}

\begin{enumerate}
 \item Le mandat des personnes commissaires est le même que celui de l’exécutif national.
\end{enumerate}

\subsection{Destitution}
\label{personnes-commissaires-destitution}

\begin{enumerate}
 \item Une personne commissaire peut être destituée par une résolution spéciale du CA.
\end{enumerate}

\subsection{Vacance}
\label{personnes-commissaires-vacance}

\begin{enumerate}
 \item Si une vacance survient, le CA, sur recommandation de l'exécutif national, peut nommer une personne commissaire par intérim jusqu'à la fin du mandat.
\end{enumerate}

\subsection{Rémunération}
\label{personnes-commissaires-remuneration}

\begin{enumerate}
 \item Les personnes commissaires siégeront sans rémunération et aucune personne commissaire ne recevra directement ou indirectement un quelconque profit de l'occupation de son poste, à condition que les dépenses raisonnables engagées dans l'exercice des fonctions de la personne commissaire soient remboursées.
\end{enumerate}
